\documentclass[UTF8,11pt]{ctexart}
\usepackage{amsmath}
\usepackage{amssymb}
\usepackage[papersize={176mm,250mm},margin=2.2cm]{geometry}
\usepackage{booktabs}
\usepackage{array}
\usepackage{enumitem}
\usepackage{xcolor}
\usepackage{listings}
\usepackage[hidelinks]{hyperref}

\definecolor{codegray}{RGB}{230,230,235}
\definecolor{codered}{RGB}{180,60,60}
\lstset{
  basicstyle=\ttfamily\small,
  backgroundcolor=\color{codegray},
  frame=single,
  framerule=0.4pt,
  rulecolor=\color{gray!60},
  columns=fullflexible,
  keepspaces=true,
  showstringspaces=false,
  breaklines=true,
  breakatwhitespace=false,
  tabsize=4,
  keywordstyle=\color{codered}\bfseries,
  commentstyle=\color{gray},
}
\newcommand{\codeinl}[1]{\texttt{#1}}
% 统一的"主句"强调命令
\newcommand{\takeaway}[1]{\medskip\par\noindent\textbf{主句：}#1\par\medskip}

\title{\bfseries 趋势跟踪与仓位管理教程\\[3pt]\large 从趋势溢价到可运行的仓位纪律（深度·自含·实证）}
\author{量化投研学习笔记}
\date{}

\begin{document}

\maketitle

\begin{abstract}
本教程把期货/CTA 的两大命题一次讲通：\textbf{方向}（趋势溢价是什么、信号怎么设计、为什么值得暴露给趋势因子）与\textbf{量}（仓位管理五级纪律）。每一级都给出\textbf{逐行推导 + 一个数字例 + 边界条件 + 一句话主句}，并用真实通达信本地数据（RB 等 8 品种）把结论钉死。
\end{abstract}

\tableofcontents
\newpage

\section*{如何使用本教程}
\addcontentsline{toc}{section}{如何使用本教程}
本教程是\textbf{递进式的}，不要跳读：
\begin{enumerate}[itemsep=2pt]
\item 第 1--3 节打概念地基：要解决的问题、趋势溢价、趋势为什么能赚钱。
\item 第 4--5 节搭骨架：方向如何映射成敞口，以及“单笔—组合—容量”的层级框架。
\item 第 6--10 节逐级给仓位纪律（每级都有推导和数字例）。
\item 第 11--13 节把组合层、容量约束与统一决策流程串起来。
\item 第 14--15 节用通达信本地真实数据给出单品种与多品种的实证。
\end{enumerate}
每节末尾的 \textbf{主句} 是可记忆的一句话；\textbf{边界} 提示该结论什么时候不成立；附录的\textbf{练习自测}可用来检验理解。

\section{要解决什么问题（问题定义）}
本教程唯一要回答的问题：
\begin{quote}
\textbf{如何把“捕捉可持续大趋势”的信念，转成一套确定规模、确定风险、能被分散放大、且能扛住反复震荡的仓位系统？}
\end{quote}
拆开是三件事：
\begin{enumerate}[itemsep=2pt]
\item \textbf{方向}：何时该长、何时该短（趋势信号）——回答“暴露给谁”。
\item \textbf{量}：每个信号下该给多大敞口（仓位）——回答“暴露多少”，并且\textbf{量决定生死}。
\item \textbf{稳健}：在大量小亏 + 偶发大赚的收益分布里，如何不把自己磨死、不错过大趋势（风控 + 分散）。
\end{enumerate}

\takeaway{本教程自始至终只回答一个问题：把“抓大趋势”的信念，转成“既能放大、又能长寿”的仓位系统。}

\section{趋势溢价：你的利润主体}
\subsection{为什么是大趋势而非技巧在赚钱}
用回归把组合收益拆开：
\begin{equation}
R_p-R_f \;=\; \underbrace{\beta_{trend}\,R_{Trend}}_{\text{趋势溢价·利润主体}}
+\; \underbrace{\alpha_{sys}}_{\text{相对基准的增值}} +\; \varepsilon .
\end{equation}
\begin{itemize}[itemsep=2pt]
\item $\beta_{trend}\,R_{Trend}$：你暴露在“趋势因子”上，系统性收获趋势溢价。\textbf{这是收入主体}。
\item $\alpha_{sys}$：相对一个朴素趋势基准的增值——来自选品、波动率目标化、尾部对冲、退场纪律。它\textbf{放大} $\beta_{trend}$ 的收获，而非独立第二利润源。
\item $\varepsilon$：单期运气，可分散、不重复。
\end{itemize}
\subsection{为什么 $\beta$ 在这里不能剥}
权益市场剥市场 $\beta$ 合理，因为指数能廉价复制它；但趋势溢价\emph{不能被廉价稳定复制}——需要系统规则、严格风控与规模容量。因此趋势 $\beta$ 是\textbf{要牢牢持有、配好“量”}的资产，不是要削掉的噪声。
\begin{quote}
第 6--10 节所有“量”的纪律，都是在决定这个 $\beta_{trend}$ 该放大多少、何时收手。\textbf{别把“高 $\beta_{trend}$ 的收益”当作 $\alpha$ 去削}。
\end{quote}
\takeaway{在期货里，趋势溢价就是那个“被付钱的 $\beta$”，是利润主体；仓位纪律的全部意义是决定这份暴露给多少、何时收。}

\subsection{为什么它也叫做「风险 Alpha」}
把对“标准基准”估出的 $\alpha$ 再拆一层：
\begin{equation}
\alpha_{measured} \;=\; \underbrace{\beta_{trend}\,R_{Trend}}_{\text{风险Alpha·趋势溢价}}
\;+\; \underbrace{\text{凸性/危机}\,\alpha}_{\text{尾部保单}}
\;+\; \underbrace{\text{信息/纪律}\,\alpha}_{\text{真本事}} \;+\; \varepsilon .
\end{equation}
趋势策略是多空、相对常规基准 $\beta\approx0$，其超额按 CAPM 定义自然是“$\alpha$”；但主要来源是\textbf{承担可定价风险挣来的补偿}：
\begin{itemize}[itemsep=2pt]
\item \textbf{信息 $\alpha$}（价值/基本面错价）：零和、容量极小、靠真本事、一挤就没；
\item \textbf{风险 $\alpha$}（趋势/动量/凸性溢价）：\emph{非零和}、可被系统化低成本“造”出来、有明确的风险画像（负偏、尾部）。
\end{itemize}
趋势跟踪绝大多数“显示 $\alpha$”属于后者，故用\textbf{风险 Alpha（另类风险溢价）}与信息 $\alpha$ 区分。含义：\textbf{它不是免费的聪明，而是有代价的承受}——应“别削、管好量、组合分散”。

\section{趋势为什么能赚钱（先讲可信度）}
\subsection{收益分布的第一个事实：低胜率高赔率右偏}
趋势系统单笔胜率常只有 $30\%$--$45\%$，但赢的时候赚得远大于亏的时候。设胜率 $p$、平均盈亏比 $b$：
\begin{equation}
\mathbb{E}[\text{单笔}] = p\cdot b - (1-p) = p(b+1)-1 .
\end{equation}
只要 $p(b+1)>1$，即便胜率较低，“对的时候赢得多”与否决定期望正负。
\subsection{为什么量价会延续（微观基础）}
趋势溢价的可持续来源至少有三：
\begin{enumerate}[itemsep=2pt]
\item \textbf{信息扩散}：重大信息并非瞬间被所有参与者消化，价格沿信息方向持续移动数日到数月；
\item \textbf{机构资金与大单驱动}：大资金建仓/减仓本身推着价格延续；
\item \textbf{行为偏差}：锚定、处置效应（人过早兑现盈利、过晚认亏）令价格在趋势上过度延续。
\end{enumerate}
\subsection{一个数字例}
设 $p=0.35$、$b=2.0$。则
$$
\mathbb{E}[\text{单笔}]=0.35\times2 -0.65 = 0.05 \ (>0).
$$
每 100 次交易约 35 次盈利、平均每次赚 $2$ 单位，65 次亏损、平均每次亏 $1$ 单位：净赚 $35\times2-65\times1=5$ 单位。
\takeaway{趋势靠“右偏分布”赚钱；仓位纪律的第一要义是让这个右偏能活到“对的时候”。}

\section{从信号到敞口：方向如何变成仓位}
\subsection{一条链：方向 $\to$ 信号 $\to$ 敞口 $\to$ 手法}
\begin{center}
\renewcommand{\arraystretch}{1.35}
\begin{tabular}{@{}l>{\raggedright\arraybackslash}p{4.6cm}@{}}
\midrule
方向 & sign(未来 $H$ 日累计收益)，$H$=时序动量窗口（如 252） \\
敞口 & 本期有意承担的 $\beta_{trend}$（大小由第 6--10 节决定） \\
换手 & 信号翻转才调仓；过高换手被成本吃掉（见第 12 节） \\
\bottomrule
\end{tabular}
\end{center}
注意“方向”和“敞口”是两层：\textbf{方向回答多大/多空，敞口回答下多少注}。
\takeaway{先有方向（sign），后有量（$w_t$）；两者都用过去信息、都可被 $\mathcal{F}_{t-1}$ 验证。}

\subsection{何时进场（进场闸门要严）}
进场的职责是\textbf{过滤掉震荡假信号、只在真趋势启动时上车}。常用触发（满足其一即可）：
\begin{center}
\setlength{\tabcolsep}{4pt}
\renewcommand{\arraystretch}{1.35}
\begin{tabular}{@{}>{\raggedright\arraybackslash}p{2.8cm}>{\raggedright\arraybackslash}p{5.0cm}>{\raggedright\arraybackslash}p{3.4cm}@{}}
\toprule
方式 & 触发条件 & 优势 / 代价 \\
\midrule
通道突破 Donchian & 收盘升破近 $N$ 日最高（多）/跌破近 $N$ 日最低（空） & 确认强；但追在突破位、成本偏高 \\
动量 sign & 过去 $H$ 日收益符号翻转（$\mathrm{sign}$ 为正才做多） & 与方向一致；翻转偏慢 \\
趋势内回撤 & 上涨中回调至均线/波动收缩企稳才进 & 价格较优；易被扫损再拉起 \\
\bottomrule
\end{tabular}
\end{center}
\textbf{同时的过滤条件}：regime 判震荡则不进（第 10 节）；用第 6 节波动率目标化控进场敞口；所有触发只用 $t-1$ 及之前信息、\textbf{无前视}。

\subsection{何时出场（盈利侧松、亏损侧紧）}
\begin{itemize}[itemsep=2pt]
\item \textbf{盈利侧要松}：不设固定止盈（那会砍掉趋势右尾），用\textbf{追踪/渠道止损}——只在趋势破坏时才走。
\item \textbf{亏损侧要紧}：\textbf{硬止损}，把单笔亏损钳在固定风险内（第 7 节）。
\item 两类结构离场：\textbf{信号反转}与\textbf{时间止损}。
\end{itemize}
\takeaway{进出场拧出“止损有上限、盈利无上限”的右偏收益分布；离场是“谁先触发谁生效”。}

\subsection{合成规则与过拟合陷阱}
真实下单把上述离场条件放进\textbf{“谁先触发谁生效”的或门}。
\begin{itemize}[itemsep=2pt]
\item \textbf{过拟合高发}：通道长度 $N$、追踪倍数 $k$、止损倍数 $s$、持有上限 $T$ 每个自由度都能造出假收益；判定标准是\textbf{样本外稳健}。
\item \textbf{换手成本}：进场过频被交易门槛吃掉——离场后需冷却期。
\item \textbf{勿设固定止盈}：那等于主动砍掉右尾。
\end{itemize}

\section{仓位管理的层级框架（本教程的总地图）}
六个决策散布在三个层级：
\begin{center}
\renewcommand{\arraystretch}{1.45}
\begin{tabular}{@{}lll@{}}
\toprule
层级 & 决策 & 对应小节 \\
\midrule
单笔（单品种） & 每手下多少手、每笔敢亏多少 & 第 5、6 节 \\
单笔（上界/概率） & 单笔最多下到多大（Kelly 上界） & 第 8 节 \\
时序（动态） & 亏损中要不要减仓、趋势不再要不要停 & 第 8、9 节 \\
组合（多品种） & 钱在品种/策略间怎么分、分散多少 & 第 11 节 \\
容量（整体） & 钱多了会怎样、成本与冲击如何 & 第 12 节 \\
总装 & 一步步串成可执行的决策链 & 第 13 节 \\
\bottomrule
\end{tabular}
\end{center}
\takeaway{五级仓位纪律不是技巧列表，而是“单笔→组合→容量”一根风险预算链上不同截面的决策。}

\section{单笔纪律 1：波动率目标化（等风险头寸）}
\subsection{做什么}
按标的\textbf{当前波动}反比确定敞口，使“每期被波动吃掉的风险”大致固定。
\subsection{为什么}
收益波动聚集（vol clustering）：大波动期也是高回撤风险期。把条件方差钉在目标值，可跨品种可比、削平回撤、不依赖精确预测。
\subsection{完整推导（逐行）}
设 $r_{t+1}$ 的条件方差（用截至 $t$ 信息估计）记为 $\sigma_t^2$。估计常用：
\begin{align}
\text{EWMA/RiskMetrics:}\quad &\hat\sigma_t^2 = (1-\lambda)\sum_{j\ge1}\lambda^{j-1}(r_{t-j}-\hat\mu_{t-j})^2,\quad \lambda\approx0.94;\\
\text{滚动已实现:}\quad &\hat\sigma_t^2 = \frac1k\sum_{j=0}^{k-1}(r_{t-j}-\bar r_t)^2,\quad k=20.
\end{align}
策略施加敞口 $w_t$，条件方差 $w_t^2\sigma_t^2$。令其等于目标波动 $\sigma_{tar}^2$，解得
\begin{equation}
\boxed{\; w_t = \frac{\sigma_{tar}}{\sigma_t},\qquad
w_t = \min\Big(\frac{\sigma_{tar}}{\hat\sigma_t},\,w_{\max}\Big)}
\end{equation}
\subsection{数字例（RB）}
RB 朴素趋势年化波动 $\approx25.2\%$。设 $\sigma_{tar}=20\%$、$w_{\max}=1.5$，则 $w_t=0.20/0.252\approx0.79$（基准八折）。
\subsection{边界}
\begin{itemize}[itemsep=2pt]
\item \textbf{滞后}：波动骤升当期仍高位加仓——接受或用更短 EWMA。
\item \textbf{双重目标化}：品种层+组合层各做一次会叠出 $\sqrt N$ 倍杠杆；组合层只做一次。
\item \textbf{单品种不显著}：第 14 节 RB 上 $t\approx0.7$；到组合层才显著（第 15 节 $t\approx2.2$）。
\end{itemize}
\takeaway{波动率目标化 $w_t=\sigma_{tar}/\sigma_t$ 把每期条件风险钉在常量——最基础的一层，也是组合层可比性的前提。}

\section{单笔纪律 2：固定风险头寸（ATR 手数）}
\subsection{做什么}
把“每笔被止损的金额”固定为账户的一个百分比，倒推手数。
\subsection{为什么}
单笔风险要可控，需回答两个相互独立的问题：\textbf{(a) 这笔总共敢亏多少}（由资金 $C$ 与风险预算 $f$ 决定）；\textbf{(b) 单手持仓被止损一次亏多少}（由止损距离与合约面值决定）。手数 $N$ 只是这两个独立量之比。用 ATR 而非价格百分比，是因为 ATR 自适应波动，让“被止损金额”稳定逼近固定的 $Cf$。
\subsection{完整推导（逐行）}
每笔愿意亏的钱：$C\,f$。单手持仓若被止损亏损：$(s\cdot\text{ATR})\times m$。故手数为两者之比、向上取整：
\begin{equation}
\boxed{\; N=\left\lceil \frac{C\,f}{(s\cdot\text{ATR})\,m}\right\rceil }
\end{equation}
\subsection{数字例}
设 $C=10^6$、$f=0.5\%$、ATR$\approx30$、$s=2$、$m=10$：$N=\lceil 5000/600\rceil=9$ 手；更保守 $s=3$：$N=6$ 手。
\subsection{边界}
\begin{itemize}[itemsep=2pt]
\item \textbf{固定风险 vs 固定手数}：固定风险随净值复利、连亏自动缩手；固定手数不复利。趋势系统的连亏连涨结构适合固定风险。
\item \textbf{手数离散}：取整造成风险跳变；小资金可能一手都下不了。
\item \textbf{ATR 失效边界}：换月跳空、涨跌停、隔夜跳空时止损价可能被击穿（gap），须配合账户级熔断。
\item \textbf{单笔下注规则}：不替代组合层分散与动态降杠杆。
\end{itemize}
\takeaway{手数 $N=\lceil Cf/(s\cdot\text{ATR}\cdot m)\rceil$ 把“敢亏多少”换成“下多少手”。}

\section{单笔纪律 3：Kelly 与分数 Kelly（上界）}
\subsection{推导}
egin{equation}
f^{*}=\frac{p(b+1)-1}{b}.
\label{eq:kelly}
\end{equation}
对趋势常用 $p\approx0.35,\,b\approx2.0$，得 $f^{*}=(0.35\times3-1)/2=0.025=2.5\%$。
\subsection{为什么绝不用满 Kelly}
\begin{enumerate}[itemsep=2pt]
\item $f^{*}$ 对 $(p,b)$ 极敏感，而他们是\textbf{估计值}；微小误差可让 $f^{*}$ 翻数倍。
\item Kelly 前提“分布已知、平稳、无限可再平衡”实盘不满足。
\item 满 Kelly 回撤路径灾难性。
\end{enumerate}
\subsection{工程结论}
用\textbf{分数 Kelly}（$\tfrac12$ 或 $\tfrac14$）作\textbf{上界/护栏}：$w_t\le\tfrac12 f^{*}$，不是目标。
\takeaway{Kelly $f^{*}=[p(b+1)-1]/b$ 是理论下界上界；趋势用它的 $\tfrac12\sim\tfrac14$ 当护栏。}

\section{动态纪律 1：基于回撤的降杠杆}
\subsection{做什么}
净值从峰值回落超阈值就下调总敞口；回升后按冷却规则恢复。
\subsection{完整推导}
\textbf{（1）回撤状态变量：} $q_t=\ln E_t$, $q_t=q_{t-1}+\ln(1+r_t)$, $\bar q_t=\max_{s\le t}q_s$, $d_t=\bar q_t-q_t$. 权益回撤 $\mathrm{DD}_t=1-E_t/\bar E_t=1-e^{-d_t}$。
\textbf{（2）降杠杆缩放 $\beta$：} $r_t=g_{t-1}\,\beta_{trend}\,\mathrm{sign}_{t-1}\,x_t$，即 $g$ 只改量不改方向。
\textbf{（3）对数增长判据：} $\mu_g=g\mu_x-\tfrac12 g^2\sigma_x^2$，$\partial\mu_g/\partial g=0\Rightarrow g=f^{*}=\mu_x/\sigma_x^2$。\textbf{过杠杆}时降仓真增利；\textbf{欠杠杆}时降仓是低点砍掉应持有的敞口。
\subsection{数字例}
阈值 $a=15\%, b=25\%$：$\mathrm{DD}\in[0.15,0.25)$ 时 $g=0.5$，$\ge0.25$ 时 $g=0$。RB 基准最大回撤 $47.4\%$，按此规则深处近乎空仓——正对应 $\beta_{trend}$ 从 0.93 落到 0.53。
\subsection{冷却/恢复（watermark）}
仅当 $q_t\ge\bar q_\tau-\Delta$ 才允许加回，避免“砍在最低点再清在最低点”。
\subsection{边界}
回撤降杠杆\textbf{只在组合层评估才有价值}：单品种负优化；组合层把回撤压到 $19.6\%$。
\takeaway{回撤降杠杆缩放 $\beta_{trend}$ 的量；是否划算取决于回撤期你相对 Kelly 是否过杠杆——因此只在组合层评估才有价值。}

\section{动态纪律 2：Regime 感知调杠杆}
\subsection{做什么}
用市场状态直接设定敞口倍率：趋势强 $\to$ $1.0$；趋势弱 $\to$ $0.6$；震荡 $\to$ $0$。
\subsection{为什么}
趋势溢价的“有/无”随时间切换。趋势强 $\to$ 满敞口；震荡 $\to$ $0$/低敞口。与第 9 节互补：regime 管“现在该不该开”，回撤降杠杆管“亏多了要不要收”。
\subsection{边界}
regime 识别带滞后与误判；单段数据训练的规则有过拟合风险。应视为“第二层风控”。
\takeaway{regime 调杠杆与回撤降杠杆构成同一根决策链的“开关”与“收手”。}

\section{组合层：等风险、相关性与单次目标化}
\subsection{为什么分散是趋势的真实放大器}
趋势溢价的\textbf{时序多样性}在跨资产上真实存在——弱信号按低相关叠加后组合夏普远超单品种（第 15 节：$0.09\to0.28$）。但趋势系统彼此高度相关，加一个趋势策略常是伪分散。
\subsection{等风险（risk parity）}
权重 $w_j^{(rp)}\propto 1/\hat\sigma_j$。若 $\rho_{jl}\approx0$，组合波动 $\approx\sigma_{tar}\sqrt{N}$。
\subsection{单次目标化原则}
组合层只做一次目标化，不要再在每个品种内重复，否则 $\sqrt N$ 放大。
\takeaway{组合层用“等风险 + 相关性去伪 + 单次目标化”把分散变成可度量的杠杆收益。}

\section{容量与成本（收紧边界）}
\subsection{参与率冲击}
\begin{equation}
\text{impact} \approx K\sqrt{p},\qquad p\le 10\%.
\end{equation}
冲 击成本在 $p$ 变大时被 $\sqrt p$ 放大。
\subsection{容量上限}
容量 $\approx$ 可容纳仓位 $\times$ 标的流动性。大资金靠分散+低换手+更大时间尺度换容量。
\subsection{边界}
收益与规模在趋势里往往负相关（规模诅咒）。
\takeaway{仓位不是越大越好：参与率 $\sqrt p$ 冲击与容量上限决定这套系统能吃多大资金。}

\section{统一决策链（把第 6--12 节串成一个流程）}
\begin{enumerate}[itemsep=2pt]
\item 方向：算 $\mathrm{sign}_t$（第 4 节）;
\item 单笔：固定风险头寸得手数 $N$（第 7 节），波动率目标化校 $\beta_{trend}$ 的量（第 6 节）;
\item 上界：$w_t\le\tfrac12 f^{*}$（第 8 节）;
\item 动态：regime 判可开（第 10 节），回撤降杠杆判收（第 9 节）;
\item 组合：等风险 + 相关性去伪 + 一次目标化（第 11 节）;
\item 容量：参与率 $\le10\%$ + 容量上限（第 12 节）;
\item 复核：净值做第 2 节回归，看 $\beta_{trend}$ vs $\alpha_{sys}$。
\end{enumerate}
\takeaway{仓位管理是按“方向→单笔→上界→动态→组合→容量→复核”串成可执行、可复核的决策链。}

\section{实证示例：单品种}
以 RB（通达信本地主力连续 L8）为例：信号 \texttt{sign(过去 H=252 日收益)}；基准=固定敞口朴素趋势；叠加1 波动率目标化；叠加2 回撤降杠杆。实测：
\begin{center}
\setlength{\tabcolsep}{4pt}
\renewcommand{\arraystretch}{1.35}
\begin{tabular}{@{}>{\raggedright\arraybackslash}p{4.2cm}>{\raggedright\arraybackslash}p{1.3cm}>{\raggedright\arraybackslash}p{1.3cm}>{\raggedright\arraybackslash}p{1.2cm}>{\raggedright\arraybackslash}p{1.6cm}>{\raggedright\arraybackslash}p{1.3cm}@{}}
\toprule
策略 & 年化 & 波动 & 夏普 & 最大回撤 & Calmar \\
\midrule
朴素趋势 & 1.9\% & 25.2\% & 0.07 & 47.4\% & 0.04 \\
+波动率目标化 & 3.1\% & 21.2\% & 0.15 & 50.3\% & 0.06 \\
+回撤降杠杆 & $-1.0\%$ & 14.1\% & $-0.07$ & 58.0\% & $-0.02$ \\
\bottomrule
\end{tabular}
\end{center}
归因（vs 基准，月频 OLS）：波动目标化 $\beta_{trend}=0.93,\alpha=+1.4\%,t=0.72$（不显著）；回撤降杠杆 $\beta_{trend}=0.53,\alpha=-2.0\%$（负优化）。

\section{实证示例：多品种等权 / 等风险}
8 品种（RB、I、HC、MA、TA、M、AU、CU），H=252，每品种 $\sigma_{tar}=20\%$，组合目标 $25\%$：
\begin{center}
\setlength{\tabcolsep}{4pt}
\renewcommand{\arraystretch}{1.35}
\begin{tabular}{@{}>{\raggedright\arraybackslash}p{4.0cm}>{\raggedright\arraybackslash}p{1.3cm}>{\raggedright\arraybackslash}p{1.3cm}>{\raggedright\arraybackslash}p{1.2cm}>{\raggedright\arraybackslash}p{1.5cm}>{\raggedright\arraybackslash}p{1.3cm}@{}}
\toprule
组合 & 年化 & 波动 & 夏普 & 最大回撤 & Calmar \\
\midrule
A 等权朴素 & 3.5\% & 12.6\% & 0.28 & 23.8\% & 0.15 \\
B 等风险波动目标 & 5.4\% & 10.4\% & \textbf{0.52} & 22.9\% & 0.23 \\
C B+组合上封 & 8.3\% & 20.0\% & 0.41 & 37.8\% & 0.22 \\
D B+回撤降 & 5.1\% & 10.2\% & 0.50 & \textbf{19.6\%} & 0.26 \\
\bottomrule
\end{tabular}
\end{center}
归因（vs A）：B $\beta_{trend}=0.88,\alpha=+2.1\%,\text{IR}=0.62,t=2.23$（显著）；C $t=1.47$；D $\beta=0.86,\alpha=+2.0\%,t=1.94$。
\begin{enumerate}[itemsep=2pt]
\item \textbf{分散是放大器}：单品种夏普 $0.09\to$ 组合 $0.28$。
\item \textbf{波动目标化 $\alpha$ 到组合层才显著}：$t=0.72\to2.23$。
\item \textbf{组合层回撤降杠杆转正}：回撤 $37.8\%\to19.6\%$、$t\approx1.94$。
\end{enumerate}

\appendix
\section*{附录 A：符号表}
\addcontentsline{toc}{section}{附录 A：符号表}
\begin{center}
\setlength{\tabcolsep}{6pt}
\renewcommand{\arraystretch}{1.4}
\begin{tabular}{@{}ll@{}}
\toprule
符号 & 含义 \\
\midrule
$r_{t+1}$ & 单期资产收益 \\
$\sigma_t,\ \hat\sigma_t$ & 条件波动、其估计 \\
$w_t$ & 仓位权重 \\
$C,\ f,\ m$ & 账户资本、单笔风险预算、每点金额 \\
$\text{ATR},\ s$ & 平均真实波幅、止损倍数 \\
$N$ & 头寸手数 \\
$p,\ b$ & 胜率、平均盈亏比 \\
$f^{*}$ & Kelly 分数 \\
$q_t,\ \bar q_t,\ d_t$ & 对数净值、峰值、对数回撤 \\
$g(d)$ & 回撤降杠杆倍率 \\
$\mathrm{DD}$ & 权益回撤 \\
$\bar\beta_{trend}$ & 降杠杆后有效趋势暴露 \\
$\rho_{jl}$ & 品种间相关系数 \\
$p_{\text{part}}$ & 参与率 \\
IC / IR / $t$ & 信息系数 / 信息比率 / $t$ 值 \\
\bottomrule
\end{tabular}
\end{center}

\section*{附录 B：练习与自测}
\addcontentsline{toc}{section}{附录 B：练习与自测}
\begin{enumerate}[itemsep=2pt]
\item 用 $p=0.4,b=1.8$ 代入 $\mathbb{E}[\text{单笔}]=p(b+1)-1$，判断期望是否为正；若为正，$f^{*}$ 是多少？
\item 若某品种 $\hat\sigma_{ann}=40\%$、$\sigma_{tar}=20\%$、$w_{max}=1.0$，$w_t$ 应为多少？若 2 个不相关品种各自目标化再等权，组合波动约多少？
\item ATR 例：$C=500{,}000$、$f=1\%$、ATR$=20$、$s=2$、$m=10$，$N$ 是多少？
\item 用 $\mu_g=g\mu_x-\tfrac12g^2\sigma_x^2$ 说明：回撤期若已欠杠杆，降杠杆会怎样？
\item 组合层单次目标化为什么能避免 $\sqrt N$ 倍放大？
\item 读第 15 节，为什么 B 的 $t=2.23$ 而 C 的 $t=1.47$？
\item 用第 13 节决策链把“单品种 $t=0.72$、组合 $t=2.23$”按 7 步走一遍。
\end{enumerate}

\section*{附录 C：参考阅读}
\addcontentsline{toc}{section}{附录 C：参考阅读}
\begin{enumerate}[itemsep=2pt]
\item Moskowitz, Ooi, Pedersen (2012). \emph{Time Series Momentum}.
\item Moreira \& Muir (2017). \emph{Volatility-Managed Portfolios}.
\item Kaminski \& Lo (2014). \emph{When Do Stop-Loss Rules Stop Losses?}
\item Sandberg, et al. \emph{Systematic Trading}（卡佛）。
\end{enumerate}

\section*{收束}
\addcontentsline{toc}{section}{收束}
\begin{quote}
\textbf{趋势跟踪赚的不是“每次对”，而是“对的时候把仓位放得足够大、足够久”。}方向决定你暴露给谁（$\beta_{trend}$），仓位决定你从这份暴露里拿走多少……所有纪律最后落到第 13 节那条可执行、可复核的决策链上。
\end{quote}

\end{document}
